\honest{Affective Death Spirals}{Eliezer Yudkowsky}{2007}

The more events you explain with your favourite theory, the truer it seems; then you doubt the evidence against it and use it to explain still more. Phlogiston ``could explain just about anything, so long as it didn't have to predict it in advance.'' \nb{The theory did forbid some outcomes. Burned metals should have lost phlogiston, and the finding that some gain weight counted against it. See the notes on ``Fake Causality.''} That is ``probably'' how people come to believe that Belgium secretly controls the US banking system.

This cycle causes much error, but ``it's nothing compared to'' the one that starts from a thought that feels good: a political system to save the world, a great leader, a tonic that cures cancer. I compare the two on no case.

The mechanism is the halo effect: a perceived positive trait raises the perception of other positive traits. That sounds to me like a fissioning uranium atom whose neutrons split other atoms. Weak feeling is subcritical. An attractive person seems more honest, which, ``perhaps,'' makes them seem more attractive, but each boost produces less than one more, and it dies out. \nb{A spiral needs this step, the raised judgment raising the first one in turn, and the post offers it with ``perhaps.'' The halo studies summarized in ``The Halo Effect'' measure one trait coloring others on one occasion; none follows judgments feeding back over time.}

With intense feeling the effect reaches everywhere. A believing Communist sees the wisdom of Marx in every hamburger, every denied promotion, every election and every newspaper article ``slanted in the wrong direction.'' Each interpretation confirms the Great Idea, and it feels good, so the believer seeks out more. \nb{This is mostly the first cycle again, a theory confirmed by every use, with pleasure added. No positive trait raises the perception of another. Karl Popper described the same experience, down to the newspaper: a Marxist found confirming evidence ``on every page,'' and ``Whatever happened always confirmed it.'' The post does not mention him.}

I offer some joke names, and settle on mine: ``affective death spiral.'' The remedy comes next.
